\begin{abstract}
基于 CLIP 的零样本异常定位通常将视觉 patch 与正常/异常文本原型进行匹配。已有文本原型可以通过学习得到，并能够形成跨类别迁移的语义判别轴；但局部异常证据还取决于当前测试图像的正常外观。对于光滑物体而言，局部纹理变化可能对应缺陷；对于粗糙材质或重复结构而言，相似响应可能属于正常外观。本文将这一缺失的图像条件化基准概括为 \clipzsad{} 中的 reference gap。一个文本侧诊断现象也指向同一问题：当 prompt 中显式加入当前图像的材质、表面纹理或结构描述时，部分定位图会发生改变，说明外观上下文会影响局部响应的解释。基于这一观察，本文提出 \method{}，一种训练无关的图像条件化正常参照估计方法。该方法从冻结 CLIP 特征中选择可靠 patch 证据，用频域响应刻画局部结构变化，并基于估计得到的正常参照重新计算异常响应。主结果覆盖五个工业基准；MVTec AD 和 VisA 上的受控实验显示了预期机制：直接融合 wavelet map 的表现较弱，全局参照低于图像条件化参照，而边界感知可靠性与保守正常参照估计在报告指标上取得最佳受控结果。

\noindent\textbf{关键词：} 零样本异常定位，CLIP，图像条件化正常参照，频域响应，patch 证据选择
\end{abstract}
